\documentclass[../../main.tex]{subfiles}
\begin{document}
\begin{oframed}
  \subfile{problem}
\end{oframed}

{\bf [解答]}

与えられた数列$\{a_k\}$について
\begin{align}
  A_n=\displaystyle\prod_{k=1}^n(1-a_k), \quad
  B_n=\displaystyle\sum_{k=1}^n\dfrac{a_k}{2^{k-1}} \label{eq:1}
\end{align}
とおいて，題意の不等式
\begin{align}
  A_n > 1- B_n \quad \left(\dfrac{1}{2}<a_k<1\right) \label{eq:2}
\end{align}
を数学的帰納法によって示す．まず以下の補題を示す．

\begin{lemma}[補題]
  $n \ge 2$に対して$2^n(1-B_n)<1$である．
\end{lemma}

\begin{proof}
  \cref{eq:2}から $\dfrac{1}{2^k}<\dfrac{a_k}{2^{k-1}}$だから，$k=1,2,\cdots,n$について足し合わせて
  \begin{align}
    & \sum_{k=1}^{n}\dfrac{1}{2^k}<\sum_{k=1}^{n}\dfrac{a_k}{2^{k-1}} \\
    & \dfrac{1}{2}\cdot\dfrac{1-(1/2)^n}{1-1/2}<B_n \\
    & 1-\left(\dfrac{1}{2}\right)^n<B_n \\
    \therefore
    & 2^n(1-B_n)<2^n\cdot\left(\dfrac{1}{2}\right)^n=1
  \end{align}
  となり補題は示された．
\end{proof}

以下帰納法で\cref{eq:2}を示す．

まず$n=2$の時，\cref{eq:2}は
\begin{align}
  (1-a_1)(1-a_2)>1-a_1-\dfrac{1}{2}a_2
\end{align}
となる．変形して，
\begin{align}
  & -a_2(1-a_1)>-\dfrac{1}{2}a_2 \\
  \iff
  & a_2\left(\dfrac{1}{2}-a_1\right)<0
\end{align}
で，これは$\dfrac{1}{2}<a_k<1$から成立する．

そこで$n=k \ (\ge 2)$での成立を仮定して，$n=k+1$について考える．
仮定から
\begin{align}
  A_k>1-B_k
\end{align}
だから，両辺に$(1-a_{k+1})\ (>0)$をかけると
\begin{align}
  A_{k+1}
  &>(1-B_k)(1-a_{k+1})\\
  &=(1-B_k)-(1-B_k)a_{k+1}\\
  &>(1-B_k)-\dfrac{a_{k+1}}{2^k} \quad (\because\text{補題})\\
  &=1-B_{k+1}
\end{align}
から，$n=k+1$でも\cref{eq:2}は成立する．

以上から数学的帰納法により\cref{eq:2}は示された．\qed

\newpage

\end{document}
